\section{Capítulo~\ref{sec:pain}. A dor}

Os sensores de dor, os dois fios, a porta na medula espinhal, o botão de volume descendente, a sensibilização e a captura da atenção foram medidos diretamente; as fórmulas da porta e da sensibilização são modelos simplificados do livro; a regra pela qual a máquina escolhe o volume do alarme é hipótese.

{\footnotesize
\setlength{\tabcolsep}{3pt}\renewcommand{\arraystretch}{1.15}
\begin{longtable}{>{\raggedright}p{0.5cm}>{\raggedright}p{4.0cm}>{\raggedright}p{5.6cm}>{\raggedright}p{2.3cm}>{\raggedright\arraybackslash}p{1.7cm}}
\toprule
N.º & Proposição & Experimento e o que foi medido & Fonte & Status \\
\midrule\endfirsthead
\toprule
N.º & Proposição & Experimento e o que foi medido & Fonte & Status \\
\midrule\endhead
1 & Limiar alto: os limiares de calor dos diferentes sensores de dor vão de $41^\circ$ a $46$--$48^\circ$ & Registro de fibras isoladas. Pele de macaco: limiar mediano dos sensores lentos (C) de $41^\circ$, dos sensores rápidos do segundo tipo, $46^\circ$, do primeiro tipo, acima de $53^\circ$. Pele humana: sensores C que respondem também à pressão, $40.7^\circ$; os que não respondem à pressão, $48^\circ$. & \cite{Treede1995heat, Weidner1999cfib} & medido \\
2 & Os sensores de dor se adaptam muito menos que os de toque e ficam mais sensíveis após uma lesão leve; o enfraquecimento após aquecimento forte foi medido num tipo & Queimadura leve da pele ($50^\circ$, $100\,\text{s}$): após $5$--$10$~min, o limiar de dor em humanos e o limiar dos sensores C no macaco caem $2$--$6^\circ$; os outros sensores da pele não mostram esse desvio. Nos sensores rápidos do segundo tipo, a resposta após aquecimento forte fica suprimida. A comparação da adaptação com os sensores de toque é uma estimativa geral, sem experimento próprio aqui. & \cite{LaMotte1982hyper, Treede1995heat} & medido \\
3 & Dois fios: rápido de $5$--$30$~m/s, lento de $0.5$--$2$~m/s; tempo de chegada $t=L/v$~\eqref{eq:paintime} & Velocidade de condução: sensores de dor rápidos do macaco, em média $15$ e $25$~m/s (dois tipos); sensores C humanos, $0.8$--$1.0$~m/s. Com $L=1$~m, isso dá dezenas de milissegundos e cerca de um segundo. & \cite{Treede1995heat, Weidner1999cfib} & medido \\
4 & A dor chega duas vezes: primeiro rápida e precisa, depois surda e longa & Tempo de reação ao aquecimento do antebraço humano: de $1100$ a $400\ms$ com o aumento da temperatura; o aquecimento forte da pele com pelos dá dor dupla, o da palma, não. A primeira dor só pode ser levada por fibras mais rápidas que $6$~m/s: pela latência, correspondem a ela os sensores rápidos do segundo tipo, que não existem na palma. A divisão de papéis (``onde'' e ``quão ruim'') é interpretação. & \cite{CampbellLaMotte1983first, Treede1995heat} & medido \\
5 & Porta: o neurônio de saída da medula recebe dor e toque, e o intermediário inibitório é excitado pelo toque (fig.~\ref{fig:gate}) & O esquema da porta foi proposto como teoria. Camundongos, marcação genética e remoção de grupos de neurônios: os neurônios excitatórios com somatostatina são necessários para a dor mecânica; os neurônios inibitórios com dinorfina impedem que a entrada dos sensores de toque chegue a eles; sem esses neurônios, um toque leve provoca dor. & \cite{MelzackWall1965, Duan2014} & medido \\
6 & O toque abafa a dor & Humanos, pulsos de laser dolorosos na mão: o toque simultâneo reduz a avaliação da dor e a capacidade de distinguir a energia dos pulsos; o efeito enfraquece com a distância entre o ponto do toque e o ponto da dor dentro de um mesmo segmento da pele. & \cite{Mancini2014touch} & medido \\
7 & Suposição da teoria da porta: uma dor forte desliga por si o intermediário inibitório, e a porta se escancara & No capítulo, marcado como suposição da teoria original. O trabalho em camundongos descreve como a porta impede o toque de entrar no canal da dor; o desligamento do intermediário pela dor não é mostrado nele. & \cite{MelzackWall1965} & hipótese \\
8 & Porta~\eqref{eq:gate}: limiar de $0.18$ sem toque, $0.86$ com toque, $0.11$ com $g=2$; com $\nu=1$ o toque reduz a saída de $1$ para $0.25$, com $\nu=2$ não tem efeito; o toque sozinho não passa (fig.~\ref{fig:gatecurves}) & O cálculo com \texttt{models/\allowbreak pain\_\allowbreak gate.py}, com os pesos do texto, reproduz todos esses números. & \texttt{models/\allowbreak pain\_\allowbreak gate.py} & modelo \\
9 & As vias descendentes do tronco mudam o ganho da porta nos dois sentidos & Rato, bulbo, registro durante aquecimento da cauda: alguns neurônios disparam antes da retirada (``on''), outros se calam (``off''); uma estimulação fraca dessa região suprime o reflexo. Revisão: os mesmos circuitos tanto facilitam quanto inibem a transmissão da dor, e os opioides agem por meio deles. & \cite{Fields1983rvm, Fields2004} & medido \\
10 & A expectativa de alívio abafa a dor pelo canal opioide & Pessoas com dor aguda: naquelas em que a expectativa de alívio reduziu a dor, o bloqueio dos receptores opioides devolveu a dor ao nível das demais; nas demais, o bloqueio não mudou a dor. & \cite{Levine1978} & medido \\
11 & O cérebro escolhe o volume do alarme pelo benefício da interrupção & Não há experimento em que a regra de escolha do volume tenha sido medida. & --- & hipótese \\
12 & Sensibilização: após uma lesão, a saída da porta cresce e o limiar cai, em parte por causa da medula; ela persiste depois que o estímulo lesivo cessa & Rato: após lesão da pele, o limiar do reflexo de flexão cai e a resposta cresce; a análise eletrofisiológica mostra que parte do aumento surge na própria medula. O estímulo lesivo provoca alterações duradouras da sensibilidade, que sobrevivem ao próprio estímulo. O capítulo não dá o tempo exato de decaimento. & \cite{Woolf1983} & medido \\
13 & A sensibilização~\eqref{eq:sensitization} é uma realimentação positiva; com um laço forte, o alarme pode travar depois da cura & Decorre do modelo~\eqref{eq:sensitization} (exercício no fim do capítulo). Não há experimento mostrando que a dor persistente após a cura seja um laço travado exatamente desse tipo. & dedução no capítulo & hipótese \\
14 & A dor é uma recompensa negativa, e o aprendizado com ela segue o erro de diferença temporal & Ressonância funcional, humanos, cadeias de prenúncios da dor: a atividade do estriado ventral e da ínsula anterior coincide com os sinais de aprendizado por diferença temporal. Macaco: parte dos neurônios \nt{DOPA} é excitada por um sopro de ar desagradável e pelo seu prenúncio, outra parte é inibida. & \cite{Seymour2004, Matsumoto2009} & medido \\
15 & A dor interrompe o trabalho em curso & Humanos, estímulo elétrico doloroso e testes sonoros: a resposta ao teste $100\ms$ após o início da dor fica bem mais lenta; após $1.5\,\text{s}$ já não há diferença em relação ao controle. Uma revisão reúne esses experimentos num modelo de captura da atenção. & \cite{Crombez1996disrupt, EcclestonCrombez1999} & medido \\
16 & A retirada se fecha na medula espinhal & Rato: neurônios das camadas profundas do corno dorsal reproduzem o padrão espacial do reflexo de retirada de músculos individuais; eles não podem ser excitados antidromicamente a partir da região cervical, ou seja, são intermediários espinhais. & \cite{Schouenborg1995wr} & medido \\
\bottomrule
\end{longtable}}
